\section{Exact lifting of modular sign extensions with a column-pair correction}
\label{sec:modular-sign-lifting}

Finite-field residue constraints can retain information that an
all-odd integer kernel test misses. They also admit exact generating
polynomial and Fourier counts. Counting through code enumerators and
character sums is classical; a foundational primary source is
MacWilliams, \emph{A Theorem on the Distribution of Weights in a
Systematic Code} (1963),
\url{https://doi.org/10.1002/j.1538-7305.1963.tb04003.x}.
Here we derive a specific lifting theorem: in a range of dimensions,
every spurious modular sign extension is a two-coordinate perturbation
of an existing row. A simple column statistic subtracts all such
spurious solutions exactly. No novelty claim is made pending a
targeted prior-art review.

Let $P$ have $r$ orthogonal sign rows of length $N$. A modular sign
extension modulo $q$ is a vector $v\in\{\pm1\}^N$ satisfying
$Pv\equiv0\pmod q$; an exact extension satisfies $Pv=0$ over the
integers. Count both kinds up to negation. Write their counts as
$M_q(P)$ and $E(P)$, and put
\[
 A(P)=\#\left\{\{a,b\}:a<b,\quad
 \left|\sum_{i=1}^r p_{ia}p_{ib}\right|=r-2\right\}.
\]
This statistic is computable in $O(rN^2)$ integer operations.

\begin{theorem}[Exact modular count with a column-pair correction]
Suppose $q\ge3$ is odd,
\[
 N=4(q+1),\qquad r>\frac{4(q+1)^2}{q^2}.
\]
Then
\[
 \boxed{E(P)=M_q(P)-A(P).}
\]
More precisely, every modular sign extension which is not exact differs
in exactly two coordinates from the sign or negative of a unique
existing row. Such nonexact classes correspond bijectively to the
column pairs counted by $A(P)$.
\end{theorem}
\begin{proof}
Let $d_i=p_i\cdot v$. Any two existing rows have Hamming distance
$N/2=2(q+1)$, which is even. Hence their sign products have the same
parity, and all $d_i$ are congruent modulo $4$: either all are zero
modulo $4$, or all are $2$ modulo $4$.
Each $d_i$ is even and divisible by odd $q$. In the second case,
$|d_i|=2q$, because the next possible magnitude $6q$ exceeds
$N=4q+4$. Bessel's inequality would then give
\[
 4q^2r=\sum_i d_i^2\le N\|v\|^2=N^2,
\]
contradicting the row-count hypothesis.
Thus every $d_i$ is divisible by $4q$, and its only possibilities
are $0,4q,-4q$. Two nonzero values would contribute
$32q^2>N^2$, again contradicting Bessel. If all vanish, the extension
is exact. Otherwise there is a unique index $i$ with
$|p_i\cdot v|=4q=N-4$. After negating $v$ if necessary, it differs
from $p_i$ in exactly two coordinates $a,b$.

For every $j\ne i$, orthogonality to $v$ and $p_i$ gives
$p_{ja}p_{ia}+p_{jb}p_{ib}=0$. Consequently the column products
$p_{ja}p_{jb}$ have the same sign for all $j\ne i$ and its opposite
sign for row $i$. Their sum has absolute value $r-2$.

Conversely, each column pair with this property has a unique exceptional
row $i$, since the hypotheses imply $r>2$. Flip those two coordinates
of $p_i$. The resulting vector has inner product $N-4=4q$ with
$p_i$ and zero with every other row, so is a nonexact modular extension.
No resulting class can come from two different rows: it would then
lie within Hamming distance two of signs of two orthogonal rows,
whose mutual distance is $N/2\ge8$, whereas the triangle inequality
would bound it by four. The flipped pair is unique once the row is
fixed. This proves the bijection and count formula.
\end{proof}

In particular, if all distinct columns have absolute correlation
strictly below $r-2$, then modular sign extensions are automatically
exact. This is a sufficient \emph{lifting} theorem, not an assertion
that modular solutions always exist.

\subsection{An exact algebraic certificate and positivity criterion}
When $q$ is prime, let $c_j\in\mathbb F_q^r$ be column $j$ of $P$.
In the integral group algebra of $\mathbb F_q^r$, form
\[
 \Phi_P=\prod_{j=1}^N(X^{c_j}+X^{-c_j}),\qquad
 C_q(P)=[X^0]\Phi_P.
\]
Expanding the product counts sign choices with zero residue, so
$C_q(P)=2M_q(P)$. Under the theorem's hypotheses,
\[
 \boxed{E(P)=\tfrac12 C_q(P)-A(P).}
\]
Thus the exact computable inequality $C_q(P)>2A(P)$ suffices for a
sign extension; equality certifies that none exists. Character
orthogonality gives the equivalent Fourier expression
\[
 C_q(P)=\frac1{q^r}\sum_{a\in\mathbb F_q^r}
          \prod_{j=1}^N2\cos\left(\frac{2\pi\,a\cdot c_j}{q}\right).
\]
The group-algebra expression permits integer computation without
floating-point trigonometry. Alternatively modular row reduction of
rank $s$ enumerates at most $2^{N-s-1}$ sign assignments to information
coordinates. The structural correction avoids an additional integer
orthogonality test for every surviving modular candidate.

For a nonvacuous correction example, take rows $0,1,2,3,4,5,6,8$ of
the order-$16$ Sylvester matrix. With $q=3$ one obtains
\[
 C_3(P)=88,\qquad M_3(P)=44,\qquad A(P)=8,\qquad E(P)=36.
\]
The verifier independently computes $C_3(P)$ by integer group-algebra
multiplication and by modular information-coordinate enumeration.

\subsection{The order-twenty-four repair benchmarks}
For $N=24$, use $q=5$. The row-count condition is simply $r\ge6$.
The following exact results concern the two previously recorded
terminal frames and the two single-row deletions of the twelve-row
frame. All counts are antipodal classes.
\[
\begin{array}{c|r|r|r|r|r}
\text{frame}&r&\max_{a<b}|c_a\cdot c_b|&M_3&M_5&E\\\hline
\text{terminal twelve-row}&12&8&72&0&0\\
\text{terminal sixteen-row}&16&4&76&0&0\\
\text{delete row zero}&11&7&72&2&2\\
\text{delete row one}&11&7&182&16&16
\end{array}
\]
In every row the column statistic $A(P)$ is zero, explaining why the
modulo-$5$ candidate sets coincide exactly with the real extension
sets. Modulo $3$ alone instead gives many false positives. In
particular the twelve-row frame passes all-odd integer feasibility
and modulo-$3$ sign feasibility, but fails modulo-$5$ sign feasibility.
This is a concrete stronger residue obstruction, not a denominator
or parity obstruction.

The sixteen-row frame has rank $16$ modulo $5$, so its nonexistence
certificate tests only $2^7=128$ information assignments with elementary
finite-field arithmetic. The twelve-row frame uses $2^{11}=2048$.
The two versus sixteen exact candidates after deletion also recover
the different repair possibilities, after testing their mutual
orthogonality. Candidate existence alone still does not ensure that
enough candidates are compatible for an augmenting exchange.

\subsection{A general moment bound outside the lifting dimension}
There is a weaker sufficient condition without the special relation
$N=4(q+1)$. For even $N$ and odd $q$, let $S_q$ be the modular sign
classes and put
\[
 T_q=|S_q|,\qquad B_q=\sum_{v\in S_q}\|Pv\|^2.
\]
A nonexact modular class has at least one nonzero inner product, an
even multiple of $q$, so its energy is at least $4q^2$. Consequently
\[
 E(P)\ge T_q-\frac{B_q}{4q^2}.
\]
In particular $B_q<4q^2T_q$ certifies an exact extension. This first
moment estimate is often weak; the theorem above succeeds by fully
classifying the possible nonzero energies and their column-pair
geometry. Neither criterion currently proves positivity uniformly
for all admissible orders.
